\documentclass[12pt,a4paper]{report}
\input{../ch/T0_preamble_standalone_En}

\title{Leitfaden zur ästhetischen Forensik bei Akkordeon-Gehäusen}
\author{Ein Leitfaden für Instrumentensammler}
\date{}

\begin{document}
	
	\maketitle
	
	\section*{Hinweis zur Methodik}
	Dieser Leitfaden konzentriert sich ausschließlich auf die äußere Gestaltung (Verdeck, Furnier, Beschläge). Bitte beachten Sie, dass die technische Einordnung (Stimmplatten, Mechanik, Stimmstöcke) ein separates Kapitel darstellt und unabhängig von der äußeren Hülle betrachtet werden muss. Ein historisches Gehäuse kann durchaus spätere technische Komponenten enthalten.
	
	\section*{1. Die Jugendstil-Ästhetik (ca. 1895–1914)}
	Der Jugendstil ist keine bloße Nachahmung, sondern eine Abstraktion der Natur. Das Ziel war die Veredelung des Alltags durch organische Formen.
	
	\begin{itemize}
		\item \textbf{Linienführung:} Dominanz der sogenannten „Peitschenhieb-Linie“. Die Formen wirken fließend, schwingend und asymmetrisch. Es gibt kaum harte Kanten oder mathematische Raster.
		\item \textbf{Motive:} Florale Inspirationen wie Ranken, Blüten und Insekten. Auch wenn die Muster symmetrisch angeordnet sind, wirken sie lebendig und „gewachsen“.
		\item \textbf{Verdeck-Charakter:} Das Verdeck fungiert als dekoratives Element, das die Fläche organisch durchbricht. Die Dekoration ist flächig und „atmend“.
	\end{itemize}
	
	\section*{2. Die Ästhetik der 1920er Jahre (Art Déco)}
	Nach 1918 wandelte sich der Zeitgeist grundlegend. Die Natur wurde durch die Maschine ersetzt; das Design wurde urban und technokratisch.
	
	\begin{itemize}
		\item \textbf{Geometrische Disziplin:} Es dominieren Rechtecke, Dreiecke, Stufenformen und der Mäander (Griechischer Schlüssel). Die Geometrie ist kein Füllmittel, sondern eine architektonische Vorgabe.
		\item \textbf{Industrielle Rhythmik:} Muster wie der Mäander oder wiederkehrende Ovale sind auf die industrielle Serienfertigung (Stanzen) ausgelegt. Sie wirken kühl, berechenbar und präzise.
		\item \textbf{Furnier-Gestaltung:} Nussbaum wird großflächig eingesetzt, nicht als Untergrund für Schnitzereien. Ein einfacher schwarz-weißer Kontrastrahmen (Fadenintarsie) unterstreicht die Geometrie des Holzes.
		\item \textbf{Beschläge:} Einsatz von kühlem Neusilber oder Nickel. Die Beschläge markieren nun die Ecken, statt sie zu „verzieren“, was die Konstruktion des Instruments betont.
	\end{itemize}
	
	\section*{3. Die Falle für Sammler: Historistische Geometrie}
	Ein häufiger Fehler ist die Verwechslung von Vorkriegs-Geometrie mit dem Art Déco.
	
	\begin{itemize}
		\item \textbf{Vorkriegszeit (Geometrischer Historismus):} Hier herrscht der Horror Vacui (Angst vor der Leere). Die Fläche ist komplett mit kleinteiligen, oft überladenen geometrischen Intarsien übersät. Jedes Element ist einzeln eingepasst. Es ist eine Demonstration handwerklicher Geduld, nicht industrieller Effizienz.
		\item \textbf{Unterscheidung:} Fragen Sie sich: Wirkt das Muster wie ein dicht gewebter, handgefertigter Teppich (Vorkrieg) oder wie ein klares, rhythmisches Skelett (1920er)? 
	\end{itemize}
	
	\section*{Fazit}
	Die Dekoration eines Akkordeons ist ein Spiegel des Zeitgeistes. Während der Jugendstil die Natur in das Instrument hineinwachsen lassen wollte, wollte der Art Déco das Instrument als technisches, präzises Objekt der Moderne definieren. Achten Sie bei floralen Motiven nicht auf das "Was" (z.B. eine Lilie), sondern auf das "Wie" (ist die Linie organisch-fließend oder statisch-geometrisch?).
	
\end{document}